\documentclass{article}
\usepackage{amsmath, amssymb, amsthm}
\usepackage{graphicx}
\usepackage{enumitem}
\usepackage{hyperref}
\usepackage{geometry}
\usepackage{algorithm}
\usepackage{algpseudocode}
\usepackage{tikz}
\usetikzlibrary{shapes,arrows,positioning}

\geometry{margin=1in}

\title{Machine Learning: Comprehensive Study Guide}
\author{Based on CS181 Course Materials}
\date{\today}

\newtheorem{definition}{Definition}
\newtheorem{example}{Example}
\newtheorem{property}{Property}

\begin{document}

\maketitle
\tableofcontents
\newpage

\section{Introduction}
This study guide provides a comprehensive overview of key machine learning concepts covered in the course. For each topic, we provide definitions, mathematical formulations, intuitions, and examples to help solidify understanding.

\section{Regression}

\subsection{Linear Regression Basics}
Linear regression aims to model the relationship between a scalar dependent variable $y$ and one or more explanatory variables $x$ by fitting a linear equation to observed data.

\subsubsection{The Bias Trick}
\begin{definition}[Bias Trick]
The bias trick is a technique where we append a constant feature $x_0 = 1$ to all of our data points, which allows us to incorporate the bias term $w_0$ into our weight vector, simplifying our equations.
\end{definition}

With the bias trick, we can rewrite:
\begin{align}
f(x, w) &= w_0 + w_1x_1 + ... + w_Dx_D\\
&= w^Tx
\end{align}
where $x$ now includes $x_0 = 1$ and $w$ includes the bias term $w_0$.

\subsubsection{Least Squares Loss Function}
\begin{definition}[Least Squares Loss]
The least squares loss function measures the squared difference between the actual values and the predicted values:
\begin{align}
L(w) = \frac{1}{2}\sum_{n=1}^{N}(y_n - w^Tx_n)^2
\end{align}
\end{definition}

The factor of $\frac{1}{2}$ is included to make the derivative cleaner.

\subsubsection{Derivation of Optimal Weights}
To find the optimal weights $w^*$ that minimize the least squares loss function, we take the derivative of the loss with respect to $w$, set it to zero, and solve:

\begin{align}
\nabla L(w) &= \sum_{n=1}^{N}(y_n - w^Tx_n)(-x_n) = 0\\
\Rightarrow \sum_{n=1}^{N}y_nx_n - \sum_{n=1}^{N}(w^Tx_n)x_n &= 0
\end{align}

In matrix form, with design matrix $X$ (where each row is a data point $x_n^T$) and target vector $y$:
\begin{align}
X^Ty - X^TXw &= 0\\
\Rightarrow w^* &= (X^TX)^{-1}X^Ty
\end{align}

This is known as the normal equation for linear regression.

\subsection{Basis Functions}
\begin{definition}[Basis Function]
A basis function $\phi(x)$ is a transformation applied to the input data to change its representation before applying linear regression. This allows linear models to capture non-linear relationships in the original feature space.
\end{definition}

The transformed linear regression model becomes:
\begin{align}
y(x,w) = w^T\phi(x)
\end{align}

Common basis functions include:
\begin{itemize}
\item Polynomial: $\phi(x) = (1, x, x^2, x^3, ...)$
\item Radial Basis Functions (RBF): $\phi(x) = \exp\left(-\frac{||x-\mu||^2}{2\sigma^2}\right)$
\item Sigmoid: $\phi(x) = \frac{1}{1+\exp(-x)}$
\end{itemize}

Basis functions allow us to fit more complex patterns while still using the machinery of linear regression.

\subsection{Regularization}
\begin{definition}[Regularization]
Regularization is a technique used to prevent overfitting by adding a penalty term to the loss function that discourages large parameter values.
\end{definition}

Regularization is used when:
\begin{itemize}
\item The model is too complex and might overfit
\item The data is noisy
\item The features are correlated (multicollinearity)
\item We want to enforce sparsity in the solution
\end{itemize}

\subsubsection{Ridge Regression (L2 Regularization)}
The loss function with L2 regularization:
\begin{align}
L(w) = \frac{1}{2}\sum_{n=1}^{N}(y_n - w^Tx_n)^2 + \frac{\lambda}{2}||w||_2^2
\end{align}

The closed-form solution for ridge regression is:
\begin{align}
w_{ridge}^* = (X^TX + \lambda I)^{-1}X^Ty
\end{align}

Ridge regression prevents any individual weight from becoming too large, providing a solution with moderate values.

\subsubsection{Lasso Regression (L1 Regularization)}
The loss function with L1 regularization:
\begin{align}
L(w) = \frac{1}{2}\sum_{n=1}^{N}(y_n - w^Tx_n)^2 + \lambda||w||_1
\end{align}

Lasso regression tends to produce sparse solutions (many weights set to exactly zero), effectively performing feature selection.

\subsubsection{Elastic Net Regularization}
Elastic Net combines L1 and L2 regularization:
\begin{align}
L(w) = \frac{1}{2}\sum_{n=1}^{N}(y_n - w^Tx_n)^2 + \lambda_1||w||_1 + \lambda_2||w||_2^2
\end{align}

It balances the sparsity-inducing property of L1 with the stability of L2 regularization.

\subsection{Bias-Variance Tradeoff}
\begin{definition}[Bias-Variance Tradeoff]
The bias-variance tradeoff describes the balance between two sources of error that prevent supervised learning algorithms from generalizing beyond their training set:
\begin{itemize}
\item High bias: Model is too simple and underfits the data
\item High variance: Model is too complex and overfits the data
\end{itemize}
\end{definition}

The mean squared error of a model can be decomposed as:
\begin{align}
\text{MSE} = \text{noise} + \text{bias}^2 + \text{variance}
\end{align}

\begin{itemize}
\item Adding more features or using more complex basis functions reduces bias but increases variance.
\item Increasing regularization increases bias but reduces variance.
\item Increasing training set size typically reduces variance without increasing bias.
\end{itemize}

\subsection{Bayesian Linear Regression}
In Bayesian linear regression, we place a prior distribution on the weights $w$:
\begin{align}
p(w) = \mathcal{N}(w|0, S_0^{-1})
\end{align}

The likelihood of the data is:
\begin{align}
p(y|X, w, \beta) = \mathcal{N}(y|Xw, \beta^{-1}I)
\end{align}

By Bayes' rule, the posterior distribution of weights is:
\begin{align}
p(w|X, y, \beta) \propto p(y|X, w, \beta)p(w)
\end{align}

With a normal prior and normal likelihood, the posterior is also normal:
\begin{align}
p(w|X, y, \beta) = \mathcal{N}(w|\mu_N, S_N^{-1})
\end{align}
where:
\begin{align}
S_N &= S_0 + \beta X^TX\\
\mu_N &= S_N^{-1}(S_0\mu_0 + \beta X^Ty)
\end{align}

The connection to regularization: the prior in Bayesian linear regression is equivalent to adding regularization. A zero-mean Gaussian prior corresponds to L2 regularization, with the prior precision controlling the regularization strength.

\subsubsection{Posterior Predictive Distribution}
The posterior predictive distribution gives us uncertainty in our predictions for new inputs $x_*$:
\begin{align}
p(y_*|x_*, X, y) = \int p(y_*|x_*, w)p(w|X, y)dw = \mathcal{N}(y_*|x_*^T\mu_N, \sigma_N^2(x_*))
\end{align}
where $\sigma_N^2(x_*) = \beta^{-1} + x_*^TS_N^{-1}x_*$.

\subsubsection{Marginal Likelihood}
The marginal likelihood (evidence) is:
\begin{align}
p(y|X) = \int p(y|X,w)p(w)dw
\end{align}
This measures how well a model explains the observed data, marginalizing over all possible parameter values.

\section{Classification}

\subsection{Loss Functions for Classification}
Classification aims to predict discrete class labels rather than continuous values. Several loss functions are used for classification:

\subsubsection{0/1 Loss}
\begin{definition}[0/1 Loss]
The 0/1 loss assigns a loss of 1 for misclassification and 0 for correct classification.
\end{definition}

While intuitive, 0/1 loss is non-differentiable and thus difficult to optimize.

\subsubsection{Hinge Loss}
\begin{definition}[Hinge Loss]
The hinge loss is defined as:
\begin{align}
L(w) = \sum_{i=1}^{N}\max(0, 1 - y_i(w^Tx_i))
\end{align}
where $y_i \in \{-1, 1\}$.
\end{definition}

Hinge loss penalizes not only misclassifications but also correct classifications with low confidence (margin). It is convex and thus easier to optimize than 0/1 loss.

\subsubsection{Logistic Loss}
\begin{definition}[Logistic Loss]
The logistic loss (cross-entropy loss) for binary classification is:
\begin{align}
L(w) = -\sum_{i=1}^{N}\{y_i\log(\hat{y}_i) + (1-y_i)\log(1-\hat{y}_i)\}
\end{align}
where $\hat{y}_i = \sigma(w^Tx_i)$ and $\sigma$ is the sigmoid function.
\end{definition}

\subsection{Gradient Descent}
\begin{definition}[Gradient Descent]
Gradient descent is an iterative optimization algorithm for finding the minimum of a function by moving in the direction of steepest descent.
\end{definition}

The update rule for gradient descent is:
\begin{align}
w^{(t+1)} = w^{(t)} - \eta\nabla L(w^{(t)})
\end{align}
where $\eta > 0$ is the learning rate.

Gradient descent works intuitively by:
\begin{itemize}
\item Computing the direction of steepest increase in loss (the gradient)
\item Moving in the opposite direction to decrease loss
\item Taking steps proportional to the learning rate $\eta$
\end{itemize}

If the learning rate is too high, the algorithm might overshoot the minimum and diverge. If it's too low, convergence will be very slow.

\subsubsection{Batch vs. Stochastic Gradient Descent}
\begin{itemize}
\item \textbf{Batch Gradient Descent}: Computes the gradient using the entire dataset at each step.
\item \textbf{Stochastic Gradient Descent (SGD)}: Computes the gradient using a single randomly selected data point at each step.
\item \textbf{Mini-batch Gradient Descent}: Computes the gradient using a small random subset of the data at each step.
\end{itemize}

SGD has several advantages:
\begin{itemize}
\item Faster updates (especially with large datasets)
\item Can escape local minima due to the noise in the gradient estimates
\item Can be used for online learning
\end{itemize}

\subsection{Perceptron Algorithm}
The perceptron is a linear binary classifier that updates weights when misclassifications occur.

The perceptron update rule for a misclassified point $(x_i, y_i)$ is:
\begin{align}
w^{(t+1)} = w^{(t)} + \eta x_i y_i
\end{align}

\begin{property}
If the data is linearly separable, the perceptron algorithm is guaranteed to converge to a solution that perfectly separates the classes in a finite number of steps.
\end{property}

\subsection{Logistic Regression}
\begin{definition}[Logistic Regression]
Logistic regression is a probabilistic discriminative model that estimates the probability of a binary class using the logistic sigmoid function:
\begin{align}
p(y=1|x;w) = \sigma(w^Tx) = \frac{1}{1+\exp(-w^Tx)}
\end{align}
\end{definition}

The decision boundary of logistic regression is linear in the feature space, defined by the equation $w^Tx = 0$.

For binary classification, the log-likelihood is:
\begin{align}
\log p(\{y_i\}_{i=1}^{N}|w) = \sum_{i=1}^{N}\{y_i\log\hat{y}_i + (1-y_i)\log(1-\hat{y}_i)\}
\end{align}

The gradient of the negative log-likelihood is:
\begin{align}
\nabla E(w) = \sum_{i=1}^{N}(\hat{y}_i - y_i)x_i
\end{align}

\subsubsection{Multi-class Logistic Regression}
For multi-class problems, we use the softmax function:
\begin{align}
p(y=k|x;W) = \frac{\exp(w_k^Tx)}{\sum_{j=1}^{K}\exp(w_j^Tx)}
\end{align}

The cross-entropy loss becomes:
\begin{align}
E(W) = -\sum_{i=1}^{N}\sum_{k=1}^{K}y_{ik}\log\hat{y}_{ik}
\end{align}

The gradient with respect to $w_j$ is:
\begin{align}
\nabla_{w_j}E(W) = \sum_{i=1}^{N}(\hat{y}_{ij} - y_{ij})x_i
\end{align}

\subsection{Generative vs. Discriminative Models}
\begin{definition}[Discriminative Model]
A discriminative model learns the conditional probability $p(y|x)$ directly, focusing on finding the decision boundary between classes.
\end{definition}

\begin{definition}[Generative Model]
A generative model learns the joint probability $p(x,y) = p(x|y)p(y)$, modeling how the data was generated from each class.
\end{definition}

\subsubsection{Generative Models for Classification}
In generative models for classification:
\begin{itemize}
\item We model the class prior $p(y)$
\item We model the class-conditional distribution $p(x|y)$
\item We use Bayes' rule to get $p(y|x) = \frac{p(x|y)p(y)}{p(x)}$
\end{itemize}

\subsubsection{Naive Bayes}
\begin{definition}[Naive Bayes]
Naive Bayes assumes that features are conditionally independent given the class:
\begin{align}
p(x|y) = \prod_{d=1}^{D}p(x_d|y)
\end{align}
\end{definition}

Despite this strong independence assumption, Naive Bayes often works well in practice. It is particularly effective when:
\begin{itemize}
\item The dimensionality is high
\item Training data is limited
\item Features are indeed relatively independent
\end{itemize}

\subsubsection{Comparing Discriminative and Generative Models}
\begin{itemize}
\item Logistic Regression (discriminative): Models $p(y|x)$ directly
\item Naive Bayes (generative): Models $p(x|y)$ and $p(y)$ to get $p(y|x)$
\end{itemize}

Generative models can:
\begin{itemize}
\item Generate new samples from the data distribution
\item Incorporate prior knowledge about the data generation process
\item Work well with missing data
\item Often converge faster with less data
\end{itemize}

Discriminative models:
\begin{itemize}
\item Often achieve better accuracy when training data is plentiful
\item Don't waste resources modeling aspects of the data irrelevant to the classification boundary
\item Can be more robust to model misspecification
\end{itemize}

\section{Neural Networks \& Model Selection}

\subsection{Neural Network Basics}
\begin{definition}[Neural Network]
A neural network is a computational model inspired by the human brain, consisting of interconnected nodes (neurons) organized in layers. Each neuron applies a non-linear transformation to a weighted sum of its inputs.
\end{definition}

\subsubsection{Activation Functions}
\begin{definition}[Activation Function]
An activation function introduces non-linearity into neural networks, allowing them to learn complex patterns.
\end{definition}

Activation functions are essential because:
\begin{itemize}
\item Without them, a multi-layer network would collapse to a single-layer linear model
\item They allow the network to approximate arbitrary functions
\item They determine the range of output values and affect gradient flow
\end{itemize}

Common activation functions include:
\begin{itemize}
\item Sigmoid: $\sigma(z) = \frac{1}{1+e^{-z}}$
\item Hyperbolic Tangent (tanh): $\text{tanh}(z) = \frac{e^z - e^{-z}}{e^z + e^{-z}}$
\item Rectified Linear Unit (ReLU): $\text{ReLU}(z) = \max(0, z)$
\end{itemize}

\subsubsection{Forward Propagation}
For a feed-forward neural network with $L$ layers:
\begin{align}
z^{(1)} &= W^{(1)}x + b^{(1)}\\
a^{(1)} &= h(z^{(1)})\\
&\vdots\\
z^{(L)} &= W^{(L)}a^{(L-1)} + b^{(L)}\\
\hat{y} &= \sigma(z^{(L)})
\end{align}
where $h$ is the activation function, $W^{(l)}$ is the weight matrix for layer $l$, and $b^{(l)}$ is the bias vector.

\subsection{Backpropagation}
\begin{definition}[Backpropagation]
Backpropagation is an algorithm for efficiently computing gradients in neural networks by propagating the error backward through the network.
\end{definition}

Backpropagation is useful because:
\begin{itemize}
\item Computing gradients directly would be computationally inefficient
\item It leverages the chain rule to reuse computed quantities
\item It scales well to deep networks
\end{itemize}

The backpropagation algorithm works as follows:
\begin{enumerate}
\item Perform a forward pass to compute activations for all layers
\item Compute the error at the output layer
\item Propagate the error backward through the network, computing gradients for each parameter
\item Update the parameters using the computed gradients
\end{enumerate}

For the output layer $L$, the error is:
\begin{align}
\delta^{(L)} = \frac{\partial E}{\partial z^{(L)}} = \hat{y} - y
\end{align}

For hidden layers $l$, the error is:
\begin{align}
\delta^{(l)} = \frac{\partial E}{\partial z^{(l)}} = ((W^{(l+1)})^T\delta^{(l+1)}) \odot h'(z^{(l)})
\end{align}
where $\odot$ is element-wise multiplication.

The gradients for weights and biases are:
\begin{align}
\frac{\partial E}{\partial W^{(l)}} &= \delta^{(l)}(a^{(l-1)})^T\\
\frac{\partial E}{\partial b^{(l)}} &= \delta^{(l)}
\end{align}

\subsection{Neural Networks for Regression and Classification}
Neural networks can be used for both regression and classification tasks by choosing appropriate:
\begin{itemize}
\item Output layer architecture
\item Activation functions
\item Loss functions
\end{itemize}

\subsubsection{Regression with Neural Networks}
For regression:
\begin{itemize}
\item Output layer: Linear activation (identity function)
\item Loss function: Mean Squared Error (MSE)
\begin{align}
E = \frac{1}{N}\sum_{i=1}^{N}(y_i - \hat{y}_i)^2
\end{align}
\end{itemize}

\subsubsection{Classification with Neural Networks}
For binary classification:
\begin{itemize}
\item Output layer: One neuron with sigmoid activation
\item Loss function: Binary Cross-Entropy
\begin{align}
E = -\frac{1}{N}\sum_{i=1}^{N}[y_i\log\hat{y}_i + (1-y_i)\log(1-\hat{y}_i)]
\end{align}
\end{itemize}

For multi-class classification:
\begin{itemize}
\item Output layer: K neurons with softmax activation
\item Loss function: Categorical Cross-Entropy
\begin{align}
E = -\frac{1}{N}\sum_{i=1}^{N}\sum_{k=1}^{K}y_{ik}\log\hat{y}_{ik}
\end{align}
\end{itemize}

\subsection{Model Selection and Regularization}
\begin{definition}[Model Selection]
Model selection is the process of choosing the best model architecture and hyperparameters for a given task.
\end{definition}

\subsubsection{Cross-Validation}
Cross-validation helps prevent overfitting by:
\begin{itemize}
\item Dividing the data into training and validation sets
\item Training the model on the training set
\item Evaluating it on the validation set
\item Selecting the model with the best validation performance
\end{itemize}

K-fold cross-validation divides the data into K subsets (folds), training on K-1 folds and validating on the remaining fold, repeating K times.

\subsubsection{Regularization in Neural Networks}
Regularization techniques for neural networks include:
\begin{itemize}
\item L1/L2 regularization: Adding penalty terms to the loss function based on the norm of the weights
\begin{align}
E_{reg} = E + \lambda\sum_{l=1}^{L}\sum_{i,j}(W_{ij}^{(l)})^2
\end{align}
\item Dropout: Randomly setting a fraction of neurons to zero during training
\item Early stopping: Stopping training when validation error starts to increase
\item Data augmentation: Creating additional training data through transformations
\end{itemize}

\section{Support Vector Machines (SVMs)}

\subsection{Maximum Margin Classification}
\begin{definition}[Maximum Margin Classifier]
A maximum margin classifier finds the decision boundary that maximizes the distance between the boundary and the closest data points from each class.
\end{definition}

Support Vector Machines (SVMs) are maximum margin classifiers that aim to find the hyperplane that maximizes this margin.

\subsection{Hard-Margin SVM}
For linearly separable data, the hard-margin SVM optimization problem is:
\begin{align}
\min_{w, w_0} &\frac{1}{2}||w||^2\\
\text{s.t. } &y_i(w^Tx_i + w_0) \geq 1, \text{ for all } i
\end{align}

The margin is given by $\frac{1}{||w||}$, so minimizing $||w||$ maximizes the margin.

\subsection{Soft-Margin SVM}
For non-linearly separable data, the soft-margin SVM introduces slack variables $\xi_i \geq 0$:
\begin{align}
\min_{w, w_0, \xi} &\frac{1}{2}||w||^2 + C\sum_{i=1}^{N}\xi_i\\
\text{s.t. } &y_i(w^Tx_i + w_0) \geq 1 - \xi_i, \text{ for all } i\\
&\xi_i \geq 0, \text{ for all } i
\end{align}

The regularization parameter $C$ controls the trade-off between maximizing the margin and minimizing the classification error:
\begin{itemize}
\item Large $C$: Focus on minimizing violations (low bias, high variance)
\item Small $C$: Focus on maximizing margin (high bias, low variance)
\end{itemize}

\subsection{Hinge Loss in SVMs}
The soft-margin SVM optimization is equivalent to minimizing:
\begin{align}
\frac{1}{2}||w||^2 + C\sum_{i=1}^{N}\max(0, 1 - y_i(w^Tx_i + w_0))
\end{align}

The second term is the hinge loss, which penalizes:
\begin{itemize}
\item Misclassified points ($y_i(w^Tx_i + w_0) < 0$)
\item Correctly classified points that are too close to the boundary ($0 < y_i(w^Tx_i + w_0) < 1$)
\end{itemize}

\subsection{Kernel Trick}
\begin{definition}[Kernel Trick]
The kernel trick allows SVMs to operate in a high-dimensional feature space without explicitly computing the transformation of the input data.
\end{definition}

Instead of computing $\phi(x_i)^T\phi(x_j)$ directly, we use a kernel function $K(x_i, x_j) = \phi(x_i)^T\phi(x_j)$.

Common kernel functions include:
\begin{itemize}
\item Linear: $K(x_i, x_j) = x_i^Tx_j$
\item Polynomial: $K(x_i, x_j) = (x_i^Tx_j + 1)^d$
\item Radial Basis Function (RBF): $K(x_i, x_j) = \exp(-\gamma||x_i - x_j||^2)$
\end{itemize}

The kernel trick is useful because:
\begin{itemize}
\item It enables efficient computation in potentially infinite-dimensional spaces
\item It allows non-linear decision boundaries in the original feature space
\item It can be applied to various algorithms beyond SVMs
\end{itemize}

\subsection{Support Vectors}
\begin{definition}[Support Vectors]
Support vectors are the data points that lie on the margin boundaries or violate the margin. They are the only points that influence the position of the decision boundary.
\end{definition}

Support vectors are important because:
\begin{itemize}
\item They determine the decision boundary (removing other points doesn't change it)
\item They are typically a small subset of the training data
\item They allow for efficient prediction on new data
\end{itemize}

\subsection{SVM vs Logistic Regression}
Both SVM and logistic regression produce linear decision boundaries but differ in how they:
\begin{itemize}
\item Measure error: SVMs use hinge loss, logistic regression uses logistic loss
\item Handle outliers: SVMs can be more robust due to the maximum margin
\item Provide probabilities: Logistic regression naturally gives probabilities, SVMs don't
\item Scale with dataset size: SVMs can be more efficient for medium-sized datasets
\end{itemize}

\section{Clustering}

\subsection{K-means Clustering}
\begin{definition}[K-means Clustering]
K-means clustering aims to partition N observations into K clusters, where each observation belongs to the cluster with the nearest mean.
\end{definition}

\subsubsection{K-means Objective}
The K-means objective is to minimize the within-cluster sum of squares:
\begin{align}
J = \sum_{k=1}^{K}\sum_{i \in C_k}||x_i - \mu_k||^2
\end{align}
where $C_k$ is the set of points in cluster $k$ and $\mu_k$ is the mean of cluster $k$.

\subsubsection{Lloyd's Algorithm}
Lloyd's algorithm is the standard approach for K-means clustering:
\begin{algorithm}
\caption{Lloyd's Algorithm for K-means}
\begin{algorithmic}[1]
\State Initialize cluster centers $\mu_1, \mu_2, \ldots, \mu_K$
\Repeat
\State Assign each point to the nearest cluster center:
      $r_{nk} = \begin{cases} 
      1 & \text{if } k = \arg\min_{k'} ||x_n - \mu_{k'}||^2 \\
      0 & \text{otherwise}
      \end{cases}$
\State Update cluster centers:
      $\mu_k = \frac{\sum_{n=1}^{N}r_{nk}x_n}{\sum_{n=1}^{N}r_{nk}}$
\Until convergence
\end{algorithmic}
\end{algorithm}

Lloyd's algorithm is guaranteed to converge to a local minimum of the objective function.

\subsubsection{K-means Extensions}
Extensions to K-means include:
\begin{itemize}
\item K-means++: A smarter initialization that improves convergence and final clustering quality
\item K-medoids: Uses actual data points as cluster centers, more robust to outliers
\end{itemize}

\subsection{Hierarchical Agglomerative Clustering (HAC)}
\begin{definition}[Hierarchical Agglomerative Clustering]
HAC builds a hierarchy of clusters by starting with each data point as its own cluster and successively merging the closest pairs of clusters.
\end{definition}

\subsubsection{HAC Algorithm}
\begin{algorithm}
\caption{Hierarchical Agglomerative Clustering}
\begin{algorithmic}[1]
\State Start with N clusters, one for each data point
\Repeat
\State Find the two closest clusters using the chosen linkage criterion
\State Merge these two clusters
\State Update distances between the new cluster and all other clusters
\Until only one cluster remains
\end{algorithmic}
\end{algorithm}

\subsubsection{Linkage Criteria}
Different linkage criteria define how the distance between clusters is measured:
\begin{itemize}
\item Single linkage (min-linkage): Minimum distance between any two points in different clusters
\begin{align}
d(C_i, C_j) = \min_{x \in C_i, y \in C_j}||x - y||
\end{align}
\item Complete linkage (max-linkage): Maximum distance between any two points in different clusters
\begin{align}
d(C_i, C_j) = \max_{x \in C_i, y \in C_j}||x - y||
\end{align}
\item Average linkage: Average distance between all pairs of points in different clusters
\begin{align}
d(C_i, C_j) = \frac{1}{|C_i||C_j|}\sum_{x \in C_i}\sum_{y \in C_j}||x - y||
\end{align}
\item Centroid linkage: Distance between cluster centroids
\begin{align}
d(C_i, C_j) = ||\mu_i - \mu_j||
\end{align}
\end{itemize}

The choice of linkage criterion affects the shape of the resulting clusters:
\begin{itemize}
\item Single linkage: Tends to produce elongated, "stringy" clusters
\item Complete linkage: Tends to produce compact, equal-diameter clusters
\item Average linkage: Strikes a balance between the two extremes
\end{itemize}

\subsubsection{Dendrogram}
A dendrogram is a tree diagram that visualizes the hierarchy of clusters in HAC. The height at which two clusters merge represents the distance between them.

\subsection{Comparing K-means and HAC}
\begin{itemize}
\item K-means:
  \begin{itemize}
  \item Requires specifying the number of clusters K in advance
  \item Produces flat partitioning (no hierarchy)
  \item Generally faster than HAC ($O(NKD)$ per iteration)
  \item Result depends on initialization
  \item Tends to find spherical clusters
  \end{itemize}
\item HAC:
  \begin{itemize}
  \item Does not require specifying the number of clusters in advance
  \item Produces a hierarchy of clusterings
  \item Slower than K-means ($O(N^2\log N)$)
  \item Deterministic (no initialization needed)
  \item Can find clusters of various shapes depending on linkage
  \end{itemize}
\end{itemize}

\section{Mixture Models and Topic Models}

\subsection{Expectation-Maximization (EM) Algorithm}
\begin{definition}[Expectation-Maximization]
The Expectation-Maximization (EM) algorithm is an iterative method for finding maximum likelihood estimates of parameters in statistical models with latent variables.
\end{definition}

The EM algorithm alternates between:
\begin{itemize}
\item E-step: Compute the expected value of the latent variables given the current parameters
\item M-step: Update the parameters to maximize the expected complete-data log-likelihood
\end{itemize}

The Evidence Lower Bound (ELBO) is a lower bound on the log-likelihood:
\begin{align}
\log p(X|\theta) \geq \mathcal{L}(q, \theta) = \mathbb{E}_{q(Z)}[\log p(X, Z|\theta)] - \mathbb{E}_{q(Z)}[\log q(Z)]
\end{align}
where $q(Z)$ is a distribution over the latent variables.

The EM algorithm:
\begin{itemize}
\item E-step: Find $q^{(t+1)} = \arg\max_q \mathcal{L}(q, \theta^{(t)})$, which gives $q^{(t+1)}(Z) = p(Z|X, \theta^{(t)})$
\item M-step: Find $\theta^{(t+1)} = \arg\max_\theta \mathcal{L}(q^{(t+1)}, \theta)$
\end{itemize}

\subsection{Mixture Models}
\begin{definition}[Mixture Model]
A mixture model is a probabilistic model that assumes the data is generated from a mixture of several probability distributions.
\end{definition}

The general mixture model assumes each data point $x_i$ is generated as follows:
\begin{itemize}
\item Sample a component $z_i \sim \text{Categorical}(\pi)$
\item Sample a data point $x_i \sim p(x|z_i, \theta)$
\end{itemize}

The marginal distribution is:
\begin{align}
p(x|\pi, \theta) = \sum_{k=1}^{K}\pi_k p(x|z=k, \theta_k)
\end{align}

\subsubsection{Gaussian Mixture Model (GMM)}
In a Gaussian Mixture Model, each component distribution is a Gaussian:
\begin{align}
p(x|\pi, \mu, \Sigma) = \sum_{k=1}^{K}\pi_k \mathcal{N}(x|\mu_k, \Sigma_k)
\end{align}

EM for GMM:
\begin{itemize}
\item E-step: Compute responsibilities
\begin{align}
\gamma_{ik} = \frac{\pi_k \mathcal{N}(x_i|\mu_k, \Sigma_k)}{\sum_{j=1}^{K}\pi_j \mathcal{N}(x_i|\mu_j, \Sigma_j)}
\end{align}
\item M-step: Update parameters
\begin{align}
\pi_k^{new} &= \frac{1}{N}\sum_{i=1}^{N}\gamma_{ik}\\
\mu_k^{new} &= \frac{\sum_{i=1}^{N}\gamma_{ik}x_i}{\sum_{i=1}^{N}\gamma_{ik}}\\
\Sigma_k^{new} &= \frac{\sum_{i=1}^{N}\gamma_{ik}(x_i - \mu_k^{new})(x_i - \mu_k^{new})^T}{\sum_{i=1}^{N}\gamma_{ik}}
\end{align}
\end{itemize}

\subsection{Directed Graphical Models (DGMs)}
\begin{definition}[Directed Graphical Model]
A directed graphical model (DGM) is a way to represent the conditional independence structure in a probability distribution using a directed graph.
\end{definition}

In a DGM:
\begin{itemize}
\item Nodes represent random variables or parameters
\item Directed edges represent conditional dependencies
\item Shaded nodes represent observed variables
\item Unshaded nodes represent latent variables or parameters
\item Plates represent repeated structures
\end{itemize}

\subsubsection{Expressing Models as DGMs}
For example, a Gaussian Mixture Model can be represented as:
\begin{itemize}
\item A node for the mixture weights $\pi$
\item A plate containing $N$ data points, each with:
  \begin{itemize}
  \item A latent variable $z_i$ (unshaded)
  \item An observed variable $x_i$ (shaded)
  \end{itemize}
\item Arrows from $\pi$ to $z_i$ and from $z_i$ to $x_i$
\end{itemize}

\section{Dimensionality Reduction}

\subsection{Principal Component Analysis (PCA)}
\begin{definition}[Principal Component Analysis]
PCA is a dimensionality reduction technique that finds a low-dimensional representation of the data that captures as much of the variance as possible.
\end{definition}

\subsubsection{Variance Preservation View}
PCA finds the directions (principal components) along which the data varies the most. Mathematically, the first principal component $w_1$ maximizes:
\begin{align}
w_1 = \arg\max_{||w||=1} \frac{1}{N}\sum_{i=1}^{N}(w^Tx_i)^2 = \arg\max_{||w||=1} w^TSw
\end{align}
where $S = \frac{1}{N}X^TX$ is the empirical covariance matrix.

The solution is the eigenvector of $S$ with the largest eigenvalue. Subsequent principal components are the eigenvectors with decreasingly smaller eigenvalues, subject to being orthogonal to all previous principal components.

\subsubsection{Reconstruction Loss View}
PCA can also be viewed as finding a low-dimensional representation that minimizes the reconstruction error when projecting back to the original space:
\begin{align}
\min_W \sum_{i=1}^{N}||x_i - WW^Tx_i||^2
\end{align}
where $W$ is a matrix with orthonormal columns.

Both views lead to the same solution: the columns of $W$ are the principal components (eigenvectors of $S$).

\subsubsection{SVD and PCA}
PCA is closely related to the Singular Value Decomposition (SVD) of the data matrix $X$:
\begin{align}
X = U\Sigma V^T
\end{align}

The principal components are the columns of $V$, and the singular values in $\Sigma$ are related to the eigenvalues of the covariance matrix by a factor of $\sqrt{N}$.

\subsection{Applications of PCA}
PCA is used for:
\begin{itemize}
\item Visualization: Reducing data to 2 or 3 dimensions for plotting
\item Noise reduction: Discarding dimensions with low variance
\item Feature extraction: Using the principal components as new features
\item Data compression: Storing a low-dimensional representation
\end{itemize}

\section{Graphical Models and Bayes Nets}

\subsection{Bayesian Networks}
\begin{definition}[Bayesian Network]
A Bayesian Network (BN) is a directed acyclic graph (DAG) where nodes represent random variables and edges represent direct dependencies.
\end{definition}

In a Bayesian Network:
\begin{itemize}
\item Each node $X_i$ is conditionally independent of its non-descendants given its parents
\item The joint probability factorizes as:
\begin{align}
p(X_1, X_2, \ldots, X_n) = \prod_{i=1}^{n}p(X_i|pa(X_i))
\end{align}
where $pa(X_i)$ are the parents of $X_i$
\end{itemize}

\subsubsection{Conditional Probability Tables (CPTs)}
For discrete variables, the conditional distributions $p(X_i|pa(X_i))$ are represented using Conditional Probability Tables (CPTs).

\subsubsection{Plate Notation}
Plate notation is used to represent repeated structures in graphical models. A plate (rectangle) around nodes indicates that those nodes are replicated, with the number of replications shown in the corner.

\subsection{D-Separation and Independence}
\begin{definition}[D-Separation]
D-separation is a criterion for determining whether sets of variables are conditionally independent given a Bayesian network structure.
\end{definition}

Two variables $X$ and $Y$ are d-separated by a set of variables $Z$ if and only if all paths between $X$ and $Y$ are blocked by $Z$.

A path is blocked if it contains:
\begin{itemize}
\item A chain $A \rightarrow B \rightarrow C$ where $B$ is in $Z$
\item A fork $A \leftarrow B \rightarrow C$ where $B$ is in $Z$
\item A collider $A \rightarrow B \leftarrow C$ where $B$ and all its descendants are not in $Z$
\end{itemize}

If $X$ and $Y$ are d-separated by $Z$, then $X$ and $Y$ are conditionally independent given $Z$.

\subsection{Variable Elimination}
\begin{definition}[Variable Elimination]
Variable elimination is an algorithm for exact inference in Bayesian networks by eliminating variables one by one from the joint distribution.
\end{definition}

For example, to compute $p(X_1|X_2=x_2)$:
\begin{align}
p(X_1|X_2=x_2) &= \frac{p(X_1, X_2=x_2)}{p(X_2=x_2)}\\
&= \frac{\sum_{X_3, \ldots, X_n}p(X_1, X_2=x_2, X_3, \ldots, X_n)}{p(X_2=x_2)}
\end{align}

The efficiency of variable elimination depends on the elimination order. A good heuristic is to eliminate leaf nodes first, as this typically results in smaller intermediate factors.

\section{Hidden Markov Models}

\subsection{HMM Structure}
\begin{definition}[Hidden Markov Model]
A Hidden Markov Model (HMM) is a statistical model where the system is assumed to be a Markov process with unobservable (hidden) states and observable emissions.
\end{definition}

An HMM is characterized by:
\begin{itemize}
\item A set of hidden states $S = \{s_1, s_2, \ldots, s_K\}$
\item A set of possible emissions $E = \{e_1, e_2, \ldots, e_M\}$
\item Initial state distribution $\pi$ where $\pi_i = p(s_1 = i)$
\item Transition probabilities $A$ where $A_{ij} = p(s_{t+1} = j|s_t = i)$
\item Emission probabilities $B$ where $B_{ik} = p(x_t = k|s_t = i)$
\end{itemize}

The joint probability of a sequence of states and emissions is:
\begin{align}
p(s_1, \ldots, s_n, x_1, \ldots, x_n) = p(s_1)\prod_{t=1}^{n-1}p(s_{t+1}|s_t)\prod_{t=1}^{n}p(x_t|s_t)
\end{align}

\subsection{Inference in HMMs}
Common inference questions in HMMs include:
\begin{itemize}
\item Filtering: $p(s_t|x_1, \ldots, x_t)$ - What is the current state given all observations up to now?
\item Prediction: $p(s_{t+k}|x_1, \ldots, x_t)$ - What will be the state in the future?
\item Smoothing: $p(s_t|x_1, \ldots, x_n)$ - What was the state at time $t$ given all observations?
\item Most likely path: $\arg\max_{s_1, \ldots, s_n} p(s_1, \ldots, s_n|x_1, \ldots, x_n)$ - What is the most likely sequence of states?
\end{itemize}

\subsection{Forward-Backward Algorithm}
The Forward-Backward algorithm computes:
\begin{itemize}
\item Forward probabilities $\alpha_t(i) = p(x_1, \ldots, x_t, s_t = i)$
\item Backward probabilities $\beta_t(i) = p(x_{t+1}, \ldots, x_n|s_t = i)$
\end{itemize}

Forward recursion:
\begin{align}
\alpha_1(i) &= \pi_i B_{i,x_1}\\
\alpha_{t+1}(j) &= B_{j,x_{t+1}}\sum_{i=1}^{K}\alpha_t(i)A_{ij}
\end{align}

Backward recursion:
\begin{align}
\beta_n(i) &= 1\\
\beta_t(i) &= \sum_{j=1}^{K}A_{ij}B_{j,x_{t+1}}\beta_{t+1}(j)
\end{align}

These can be used to compute:
\begin{align}
p(s_t = i|x_1, \ldots, x_n) &= \frac{\alpha_t(i)\beta_t(i)}{p(x_1, \ldots, x_n)}\\
p(s_t = i, s_{t+1} = j|x_1, \ldots, x_n) &= \frac{\alpha_t(i)A_{ij}B_{j,x_{t+1}}\beta_{t+1}(j)}{p(x_1, \ldots, x_n)}
\end{align}

\subsection{Viterbi Algorithm}
The Viterbi algorithm finds the most likely sequence of hidden states:
\begin{align}
\arg\max_{s_1, \ldots, s_n} p(s_1, \ldots, s_n|x_1, \ldots, x_n)
\end{align}

It uses dynamic programming to compute:
\begin{align}
\delta_1(i) &= \pi_i B_{i,x_1}\\
\delta_t(j) &= \max_{i}[\delta_{t-1}(i)A_{ij}]B_{j,x_t}
\end{align}
and keeps track of the argmax at each step to reconstruct the optimal path.

\subsection{Learning in HMMs}
The parameters of an HMM can be learned using the EM algorithm (Baum-Welch algorithm):
\begin{itemize}
\item E-step: Compute expected counts of transitions and emissions using the Forward-Backward algorithm
\item M-step: Update parameters to maximize the expected complete-data log-likelihood
\end{itemize}

\section{Markov Decision Processes}

\subsection{MDP Framework}
\begin{definition}[Markov Decision Process]
A Markov Decision Process (MDP) is a mathematical framework for modeling decision-making in situations where outcomes are partly random and partly under the control of a decision-maker.
\end{definition}

An MDP is defined by:
\begin{itemize}
\item A set of states $S$
\item A set of actions $A$
\item Transition probabilities $P(s'|s, a)$ - the probability of transitioning to state $s'$ when taking action $a$ in state $s$
\item Reward function $R(s, a)$ - the immediate reward received after taking action $a$ in state $s$
\item Discount factor $\gamma \in [0, 1]$ - how much future rewards are valued compared to immediate rewards
\end{itemize}

\subsection{Value Functions}
\begin{definition}[Value Function]
The value function $V^\pi(s)$ is the expected return starting from state $s$ and following policy $\pi$.
\end{definition}

For a policy $\pi$, the value function is:
\begin{align}
V^\pi(s) = \mathbb{E}_\pi\left[\sum_{t=0}^{\infty}\gamma^t R(s_t, a_t) \mid s_0 = s\right]
\end{align}

\begin{definition}[Q-Function]
The Q-function $Q^\pi(s, a)$ is the expected return starting from state $s$, taking action $a$, and then following policy $\pi$.
\end{definition}

The Q-function is:
\begin{align}
Q^\pi(s, a) = R(s, a) + \gamma\sum_{s'}P(s'|s, a)V^\pi(s')
\end{align}

The optimal value function $V^*(s)$ is the maximum value achievable from state $s$:
\begin{align}
V^*(s) = \max_\pi V^\pi(s)
\end{align}

The optimal Q-function $Q^*(s, a)$ is:
\begin{align}
Q^*(s, a) = R(s, a) + \gamma\sum_{s'}P(s'|s, a)V^*(s')
\end{align}

\subsection{Bellman Equations}
The Bellman equations are recursive relationships for value functions:

\begin{definition}[Bellman Consistency Equation]
For a given policy $\pi$:
\begin{align}
V^\pi(s) = \sum_a \pi(a|s)\left[R(s, a) + \gamma\sum_{s'}P(s'|s, a)V^\pi(s')\right]
\end{align}
\end{definition}

\begin{definition}[Bellman Optimality Equation]
For the optimal policy:
\begin{align}
V^*(s) = \max_a\left[R(s, a) + \gamma\sum_{s'}P(s'|s, a)V^*(s')\right]
\end{align}
\end{definition}

These equations form the basis for planning and learning algorithms in MDPs.

\subsection{Planning in MDPs}
\subsubsection{Finite Horizon Planning}
For a finite time horizon $T$, we compute:
\begin{align}
V_t(s) = \max_a\left[R(s, a) + \sum_{s'}P(s'|s, a)V_{t+1}(s')\right]
\end{align}
starting from $V_T(s) = \max_a R(s, a)$ and working backwards.

\subsubsection{Value Iteration}
Value iteration is an algorithm for infinite horizon planning:
\begin{algorithm}
\caption{Value Iteration}
\begin{algorithmic}[1]
\State Initialize $V(s) = 0$ for all $s \in S$
\Repeat
\State $\Delta \leftarrow 0$
\For{each $s \in S$}
\State $v \leftarrow V(s)$
\State $V(s) \leftarrow \max_a\left[R(s, a) + \gamma\sum_{s'}P(s'|s, a)V(s')\right]$
\State $\Delta \leftarrow \max(\Delta, |v - V(s)|)$
\EndFor
\Until{$\Delta < \theta$ (a small threshold)}
\end{algorithmic}
\end{algorithm}

Value iteration converges to the optimal value function $V^*$.

\subsubsection{Policy Iteration}
Policy iteration alternates between policy evaluation and policy improvement:
\begin{algorithm}
\caption{Policy Iteration}
\begin{algorithmic}[1]
\State Initialize $\pi$ arbitrarily
\Repeat
\State \textbf{Policy Evaluation:} Compute $V^\pi(s)$ for all $s \in S$
\State \textbf{Policy Improvement:} Update $\pi(s) \leftarrow \arg\max_a\left[R(s, a) + \gamma\sum_{s'}P(s'|s, a)V^\pi(s')\right]$
\Until{policy stable}
\end{algorithmic}
\end{algorithm}

Policy iteration converges in fewer iterations than value iteration but may require more computation per iteration.

\section{Reinforcement Learning}

\subsection{Reinforcement Learning vs. Planning}
\begin{definition}[Reinforcement Learning]
Reinforcement learning is a type of machine learning where an agent learns to make decisions by interacting with an environment and receiving feedback in the form of rewards.
\end{definition}

The key difference between RL and planning is that in RL, the agent:
\begin{itemize}
\item Does not know the transition probabilities $P(s'|s, a)$
\item Does not know the reward function $R(s, a)$
\item Must learn from experience through trial and error
\end{itemize}

\subsection{Exploration vs. Exploitation}
A fundamental challenge in RL is the exploration-exploitation tradeoff:
\begin{itemize}
\item Exploration: Trying new actions to discover their effects
\item Exploitation: Using known information to maximize rewards
\end{itemize}

Strategies for balancing exploration and exploitation include:
\begin{itemize}
\item $\epsilon$-greedy: Take the best known action with probability $1-\epsilon$, and a random action with probability $\epsilon$
\item Softmax: Choose actions with probability proportional to their estimated values
\item Upper Confidence Bound (UCB): Choose actions based on both their estimated value and the uncertainty in that estimate
\end{itemize}

\subsection{Model-Free vs. Model-Based RL}
\subsubsection{Model-Based RL}
In model-based RL, the agent:
\begin{itemize}
\item Learns a model of the environment (transition probabilities and rewards)
\item Uses the model for planning (e.g., with value iteration or policy iteration)
\end{itemize}

\subsubsection{Model-Free RL}
In model-free RL, the agent:
\begin{itemize}
\item Learns value functions or policies directly from experience
\item Does not explicitly model the environment
\end{itemize}

\subsection{Q-Learning}
\begin{definition}[Q-Learning]
Q-Learning is a model-free RL algorithm that learns the optimal Q-function directly from experience.
\end{definition}

The Q-learning update rule is:
\begin{align}
Q(s_t, a_t) \leftarrow Q(s_t, a_t) + \alpha\left[r_t + \gamma\max_{a'}Q(s_{t+1}, a') - Q(s_t, a_t)\right]
\end{align}
where $\alpha$ is the learning rate.

Q-learning is an off-policy algorithm because it learns about the optimal policy regardless of the policy being followed.

\subsection{SARSA}
\begin{definition}[SARSA]
SARSA (State-Action-Reward-State-Action) is a model-free RL algorithm that learns the Q-function for the policy being followed.
\end{definition}

The SARSA update rule is:
\begin{align}
Q(s_t, a_t) \leftarrow Q(s_t, a_t) + \alpha\left[r_t + \gamma Q(s_{t+1}, a_{t+1}) - Q(s_t, a_t)\right]
\end{align}

SARSA is an on-policy algorithm because it learns about the policy being followed.

\subsubsection{Differences Between Q-Learning and SARSA}
The key differences between Q-learning and SARSA are:
\begin{itemize}
\item Q-learning uses the maximum Q-value for the next state ($\max_{a'}Q(s_{t+1}, a')$)
\item SARSA uses the Q-value of the actual next action taken ($Q(s_{t+1}, a_{t+1})$)
\item Q-learning converges to the optimal policy regardless of the exploration strategy
\item SARSA converges to the optimal policy only if exploration is gradually reduced
\item SARSA can be more conservative in risky environments
\end{itemize}

\section{Conclusion}
This comprehensive study guide has covered the key concepts and techniques in machine learning, from basic regression and classification to advanced topics like reinforcement learning. Understanding these fundamentals provides a solid foundation for applying machine learning to real-world problems and exploring more advanced topics in the field.

\end{document}